\section{Introduction}\label{sec:intro}

\subsection{Background}\label{sec:intro-background}

Fourier-based integer multiplication splits two $n$-bit nonnegative integers $A$ and $B$ into $N = \lceil n / b \rceil$ radix-$2^b$ digits of $b = \Theta(\log n)$ bits. It chooses a transform length $R_0 = \Theta(n / \log n)$ with $R_0 \ge 2N$ and evaluates a cyclic convolution of length $R_0$ at fixed-point precision $p_0 = \Theta(\log n)$ bits. A one-dimensional radix-2 discrete Fourier transform (DFT) of length $R_0$ performs $\frac{1}{2} R_0 \log_2 R_0$ complex butterfly multiplications on $p_0$-bit operands. If each $p_0$-bit product costs $\mathsf{M}(p_0)$ steps, these twiddle multiplications alone cost $\Theta(R_0 \log R_0 \cdot \mathsf{M}(p_0)) = \Theta(n \cdot \mathsf{M}(\log n))$ steps, an extra factor of $\Theta(\log \log n)$ over $n \log n$ when $\mathsf{M}(p) = \Theta(p \log p)$~\cite{HvdH21}. Sch\"onhage and Strassen~\cite{SS71} achieved $O(n \log n \log \log n)$ steps on a multitape Turing machine by working in Fermat quotient rings $\Z / (2^{2^k}+1)\Z$. F\"urer~\cite{Furer07,Furer09} reduced the $\log \log n$ factor to $2^{O(\log^* n)}$ using synthetic roots in cyclotomic rings, and Harvey, van der Hoeven, and Lecerf~\cite{HvdHL16} sharpened that bound to $O(n \log n \cdot 8^{\log^* n})$ unconditionally.

A classical algebraic tool for eliminating butterfly multiplications is the polynomial transform of Nussbaumer and Quandalle~\cite{NQ79,Nussbaumer80,Nussbaumer82}. In the negacyclic quotient ring $\Rr_r := \C[y]/(y^r + 1)$ where $r$ is a power of two, the indeterminate $y$ satisfies $y^r \equiv -1$ and has multiplicative order $2r$. Multiplying an element of $\Rr_r$ by any power of $y$ is a signed cyclic shift of its $r$ complex coefficients and requires zero scalar multiplications. Combined with negacyclic-to-cyclic weighting~\cite{SS71,CF94}, a polynomial transform of length $t_1$ over $\Rr_r$ ($t_1 \mid r$) reduces a two-dimensional transform $F_{t_1, r}$ to $t_1$ one-dimensional transforms $F_r$ and an odd-stride permutation~\cite{NQ79,Nussbaumer82}. Integer multiplication, however, requires a one-dimensional transform $F_{t_1 r}$, which differs from $F_{t_1, r}$ by the Cooley--Tukey twiddle factors $e^{2\pi i j_1 u_2 / (t_1 r)}$.

Harvey and van der Hoeven~\cite{HvdH21} achieved $O(n \log n)$ steps by bridging one-dimensional and power-of-two multidimensional transforms through Chinese remaindering. Their algorithm routes the length-$R$ transform through a $d$-dimensional Good--Thomas transform of pairwise coprime lengths and resamples each axis to a power-of-two length by Gaussian convolution, a step related to nonuniform Fourier transform (NUFFT) gridding~\cite{DR93,PST01,GL04,BMK19}. Like the present paper, it evaluates the resulting power-of-two multidimensional transforms by Nussbaumer's polynomial transforms~\cite{NQ79,Nussbaumer80,Nussbaumer82}. The proof of their main theorem takes $d = 1729$~\cite{HvdH21}. Because the source length and target power-of-two length on each axis differ, the resampling offset varies across every coordinate index and expands the grid at each recursive step. A recent preprint~\cite{OAI26} gives an $O(n(\log n)^{1-\kappa})$ algorithm on a multitape Turing machine by combining a Harvey--van der Hoeven-style architecture with finite linear network primitives; the present paper does not use it.

\subsection{Main results}\label{sec:intro-results}

We show that $O(n \log n)$ integer multiplication requires neither Chinese remaindering nor resampling between grids of different sizes. Instead, we relate $F_{t_1 r}$ and $F_{t_1, r}$ directly on a single two-dimensional power-of-two grid $R = t_1 r$. Throughout, an $n$-bit integer means a nonnegative integer in $[0, 2^n)$ (signed $n$-bit integers reduce immediately in $O(n)$ steps).

\begin{maintheorem}\label{thm:main}
There is a deterministic multitape Turing machine and a constant $C > 0$ such that, for every $n \ge 2$ and any two $n$-bit integers $A$ and $B$, given the binary representations of $A$ and $B$ on the input tape, the machine writes the $2n$-bit binary product on the output tape in at most $C n \log_2 n$ steps.
\end{maintheorem}

The machine of Theorem~\ref{thm:main} has the following structural properties. Every Fourier transform it performs has power-of-two length, and each top-level transform of length $R_0 = \Theta(n / \log n)$ recurses on $t_{1,\ell} \times r_\ell$ sub-grids of the same total size ($R_\ell = t_{1,\ell} r_\ell$) with no grid expansion, no Chinese remainder map, no prime transform length, and no resampling between grids of different sizes. Inner products of $O((\log n)^2)$-bit integers are delegated to the same algorithm one level down, backed at the bottom stage by a one-dimensional radix-2 Fourier multiplier using schoolbook arithmetic.

\begin{maintheorem}[machine-checked]\label{thm:lean}
A fixed stack machine \texttt{tmIntMul} (Mathlib's \texttt{Turing.FinTM2}: eight binary stacks, $32$ register states, a fixed finite control) satisfies the following, verified in Lean~4~\cite{dMU21,Mathlib20}. For every $n \ge 2$ and all $A, B < 2^n$, started with the $2n$ input bits it halts with the $2n$ bits of $A \cdot B$ on its output stack within $C \cdot n \lfloor \log_2 n \rfloor$ steps, for a constant $C$ (the formal statement is existential; its proof exhibits the explicit witness \texttt{C\_tm\_total} $\approx 5\cdot 10^{92}$, Section~\ref{sec:lean}, \texttt{IntMul.\allowbreak tmIntMul\_\allowbreak evalsToInTime\_\allowbreak n\_\allowbreak log\_\allowbreak n}). Every step executes $O(1)$ elementary stack operations (\texttt{IntMul.\allowbreak tmIntMul\_\allowbreak elementary\_\allowbreak stack\_\allowbreak ops\_\allowbreak n\_\allowbreak log\_\allowbreak n}). The proof uses no axioms beyond \texttt{propext}, \texttt{Classical.choice}, and \texttt{Quot.sound}.
\end{maintheorem}

\begin{table}[ht]
\centering
\small
\caption{Comparison of transform components in~\cite{HvdH21} and in this paper.}
\label{tab:components}
\setlength{\tabcolsep}{6pt}
\begin{tabular}{@{}>{\raggedright}p{0.23\textwidth}>{\raggedright}p{0.35\textwidth}>{\raggedright\arraybackslash}p{0.36\textwidth}@{}}
\toprule
Component & \cite{HvdH21} & This paper \\
\midrule
CRT map & Good--Thomas embedding onto coprime axes & None \\\addlinespace[2pt]
Coprime lengths & Pairwise coprime axis lengths & None (powers of two $t_1, r$ only) \\\addlinespace[2pt]
Resampling & Gaussian resampling between grids of different sizes & Same-grid row spreading by windowed sinc-Gaussian kernel \\\addlinespace[2pt]
Multidimensional transform & Nussbaumer polynomial transforms & Nussbaumer polynomial transform over $\C[y]/(y^r + 1)$ (one two-dimensional level) \\\addlinespace[2pt]
Dimension & High-dimensional ($d = 1729$ in main proof) & Two-dimensional ($t_1 \times r$ grid, $t_1 r = R$) \\\addlinespace[2pt]
Working precision & $O(\log n)$ bits, fixed per transform; resampling conditioned by oversampling the target grid & $O(\log n)$ bits at every level, additive schedule $p_{\ell+1} = p_\ell + 2K_\ell + 12$; two-term solvers conditioned by at most $2.36\,t_1$, no oversampling \\
\bottomrule
\end{tabular}
\end{table}

Table~\ref{tab:components} summarizes what the reduction removes relative to~\cite{HvdH21}. The tools are classical; the contribution is combining a same-grid fractional-shift Poisson bridge with a windowed sinc-Gaussian kernel, uniform two-term cyclic and Nyquist deconvolution, an additive precision schedule, and a machine-checked $O(n \log n)$ bound for an explicit eight-stack machine. Specifically, the polynomial transform over $\C[y]/(y^r + 1)$ and its reduction of $F_{t_1, r}$ to $t_1$ calls of $F_r$ with the odd-stride permutation are due to Nussbaumer and Quandalle~\cite{NQ79,Nussbaumer80,Nussbaumer82}, and the negacyclic twist is classical~\cite{SS71,CF94}. Kernel spreading analyzed by Poisson summation comes from NUFFT gridding~\cite{DR93,PST01,GL04,BMK19}, while row spreading uses Kronecker substitution~\cite{Schonhage82,Harvey09} with overlap-add alongside the classical two-call real convolution identity. Our written proof imports no external algorithmic subroutine as a black box, and its analytic inputs are standard real-analysis facts: Poisson summation for Schwartz functions, the convolution theorem and Parseval's identity, and elementary series and Gaussian tail estimates.

\subsection{Comparison with related work}\label{sec:intro-comparison}

Harvey and van der Hoeven~\cite{HvdH21} route a length-$R$ transform through a $d$-dimensional transform of pairwise coprime lengths via Chinese remaindering and resample each axis onto a larger power-of-two length by Gaussian convolution. Their recursion uses a large fixed dimension. Our algorithm works on a single two-dimensional power-of-two grid $R_\ell = t_{1,\ell} r_\ell$ at every recursion depth $\ell$. The $R_0 / R_\ell$ active subproblems at depth $\ell$ partition the original $R_0$ complex entries with zero grid expansion, while the flat passband of the windowed sinc-Gaussian kernel keeps the working precision $O(\log n)$ across all $O(\log \log n)$ levels.

A companion manuscript~\cite{CLW26} generalizes the same-grid fractional-shift bridge to $d$ dimensions with a complex half-bin phase that removes the Nyquist special case, and combines it with the finite linear network primitives of~\cite{OAI26}; it claims an $o(n \log n)$ bound on a deterministic multitape Turing machine. Nothing in the present paper depends on~\cite{CLW26}. Its main theorem treats the inner multiplier on $O((\log n)^2)$-bit operands as a black box; its fully CRT-free instantiation and its machine-checked proof use (a formalization of) the present two-dimensional multiplier as that inner multiplier. The present paper is the self-contained two-dimensional, real-kernel instance: the one that is executed up to $10^9$ bits in C/GMP and whose complexity bound is machine-checked as a statement about one concrete eight-stack machine.

\subsection{Verification and executability}\label{sec:intro-verification}

We formalize Theorem~\ref{thm:lean} in Lean~4~\cite{dMU21} on Mathlib~\cite{Mathlib20} in the library \texttt{IntMul}~\cite{LeanMS26}, comprising $27$ files and $82{,}377$ lines (Section~\ref{sec:lean}). The development defines the explicit eight-stack binary machine \texttt{tmIntMul} and proves its bit-exactness and $O(n \log n)$ step and elementary-operation bounds for all $n \ge 2$. The analytic and fixed-point bounds on which the formal proof rests (Poisson identity, passband, stopband and conditioning bounds, spreading, multi-level accuracy, root recovery) are proved inside the same library; Table~\ref{tab:lean-map} lists the correspondence we established and the places where the formal constants are looser than the printed ones. Compile-time axiom checks verify that the top-level theorems use no axioms beyond \texttt{propext}, \texttt{Classical.choice}, and \texttt{Quot.sound}.

We also evaluate the C/GMP implementation \texttt{intmul\_fastdft} (\texttt{algo=1}) of the algorithm of this paper, with the windowed sinc-Gaussian kernel and the additive schedule and inner products by GMP~\cite{CGMP26} (Section~\ref{sec:empirical}). Single-threaded on one Zen~3 core, the C/GMP program multiplies two $10^8$-bit integers in $92.5\,\mathrm{s}$ excluding $5.5\,\mathrm{s}$ of one-time table setup ($98.0\,\mathrm{s}$ including it), about $153$ times slower than GMP's \texttt{mpz\_mul} ($0.606\,\mathrm{s}$) on the same operands ($153$ to $323$ times across the four measured sizes, $162$ to $326$ times including table setup). It agrees bit for bit with GMP at every size tested (one random operand pair per size), up to $10^9$ bits.

{\sloppy\paragraph{Use of AI tools.}
This work was developed with substantial assistance from AI agents (Gemini~4 and other AI models), working under the direction of the authors. That assistance spanned the algorithm, the proofs in this paper, the C/GMP implementation and its benchmark harness, and the Lean~4 formalization. The authors have checked the written proofs in this paper. No author has checked the formal Lean proof line by line; the correctness of the formal development rests on verification by Lean's kernel.

The authors have checked the formal statements of the machine, its uniformity, and the step bound of Theorem~\ref{thm:lean} in \texttt{IntMul/TMEndToEnd.lean} and the encoding conventions in \texttt{IntMul/Defs.lean}. The further definitions and the elementary-operation theorem that fix the meaning of Theorem~\ref{thm:lean} (the \texttt{IsUniformFinTM2} predicate and the \texttt{CoeDep} instance in \texttt{IntMul/TMPrimitives.lean}, and \texttt{stmtMaxElemOps}, \texttt{FinTM2BoundedElemOps}, and \texttt{tmIntMul\_\allowbreak elementary\_\allowbreak stack\_\allowbreak ops\_\allowbreak n\_\allowbreak log\_\allowbreak n} in \texttt{IntMul.lean}) are quoted verbatim in Section~\ref{sec:lean} so that readers can inspect them. Agreement among AI systems was used only as a working heuristic during exploration, never as mathematical evidence. Every theorem and lemma stated in this paper is established by a complete written proof, an explicit citation to the literature, or a kernel-checked formal theorem. Statements about the companion manuscript report its claims, standard facts about models of computation (such as the constant-factor simulation between stack machines and multitape Turing machines) are stated as such, and the measurements of Section~\ref{sec:empirical} are read from the logs of the runs described there.\par}

\subsection{Organization}\label{sec:intro-org}

Section~\ref{sec:overview} gives an informal technical overview of the algorithm, the kernel, and the complexity accounting. Section~\ref{sec:prelim} fixes notation, the multitape Turing machine model, fixed-point grids, elementary-function tables, tape routing, and the polynomial transform. Section~\ref{sec:bridge} presents the Nussbaumer--Quandalle polynomial-transform reduction and proves the fractional-shift Poisson identity. Section~\ref{sec:deconv} analyzes the windowed sinc-Gaussian passband, the two-term cyclic deconvolution, the Nyquist alternating-sum identity, and the finite-precision column solvers. Section~\ref{sec:spread-and-error} proves the Kronecker overlap-add spreading bound and the inductive forward-error closure under the additive precision schedule, and Section~\ref{sec:main-proof} gives the complete algorithm specification and proves Theorem~\ref{thm:main}. Section~\ref{sec:lean} describes the Lean~4 formalization (Theorem~\ref{thm:lean}), Section~\ref{sec:empirical} reports the C/GMP measurements, Section~\ref{sec:discussion} discusses the remaining $\Theta(n \log n)$ term, and Appendices~\ref{app:elementary}--\ref{app:deferred} contain the deferred proofs for table evaluation, breadth-first tape scheduling, tape routing, solver rounding, and filter norms.
